% Présentation CODIR · 10 minutes + 5 minutes de questions · 6 diapositives
% Moteur : XeLaTeX. Les chiffres viennent de donnees/macros.tex (générés).
\documentclass[aspectratio=169,10pt]{beamer}
\usepackage{iftex}
\ifPDFTeX\usepackage[T1]{fontenc}\usepackage[utf8]{inputenc}\fi
\usepackage{libertinus}
\usepackage[french]{babel}
\frenchsetup{AutoSpacePunctuation=false}
\usepackage{tikz,booktabs,tabularx,array}
% Informations du groupe et de la page de garde : modifier ici, elles sont reprises partout.
\newcommand{\GroupeNom}{Équipe Team, Marc, Matteo, Bernard}
\newcommand{\Membres}{Team, Marc, Matteo et Bernard}
\newcommand{\MembresListe}{Team\\ Marc\\ Matteo\\ Bernard}
\newcommand{\Formation}{M2 SSIR}
\newcommand{\Universite}{Toulouse Paul Sabatier}
\newcommand{\Encadrant}{} % laisser vide pour masquer le bloc « Encadré par »
\newcommand{\AnneeUniv}{2026-2027}
\newcommand{\DateDepot}{10 novembre 2026}
\newcommand{\DateReference}{12 octobre 2026} % date de référence pédagogique du contrôle (T06)

% Fichier généré par scripts/generer.py depuis scripts/donnees.py. Ne pas éditer à la main.
\newcommand\Cible{4}
\newcommand\NiveauCalcule{4}
\newcommand\SommeMaxima{10}
\newcommand\Ajustement{0}
\newcommand\MaxConseq{3}
\newcommand\MaxSensib{3}
\newcommand\MaxMenaces{2}
\newcommand\MaxVuln{2}
\newcommand\Plafond{120\,000\,€}
\newcommand\TotalRetenu{116\,000\,€}
\newcommand\Marge{4\,000\,€}
\newcommand\qnote[1]{\csname qnote#1\endcsname}
\newcommand\peff[1]{\csname peff#1\endcsname}
\newcommand\rimp[1]{\csname rimp#1\endcsname}
\newcommand\rvra[1]{\csname rvra#1\endcsname}
\newcommand\rscore[1]{\csname rscore#1\endcsname}
\newcommand\acout[1]{\csname acout#1\endcsname}
\expandafter\def\csname qnoteQ01\endcsname{3}
\expandafter\def\csname qnoteQ02\endcsname{2}
\expandafter\def\csname qnoteQ03\endcsname{3}
\expandafter\def\csname qnoteQ04\endcsname{2}
\expandafter\def\csname qnoteQ05\endcsname{2}
\expandafter\def\csname qnoteQ06\endcsname{3}
\expandafter\def\csname qnoteQ07\endcsname{2}
\expandafter\def\csname qnoteQ08\endcsname{2}
\expandafter\def\csname qnoteQ09\endcsname{2}
\expandafter\def\csname qnoteQ10\endcsname{2}
\expandafter\def\csname qnoteQ11\endcsname{2}
\expandafter\def\csname qnoteQ12\endcsname{2}
\expandafter\def\csname peffStratégie\endcsname{1}
\expandafter\def\csname peffRisques\endcsname{1}
\expandafter\def\csname peffRègles\endcsname{2}
\expandafter\def\csname peffSupervision\endcsname{1}
\expandafter\def\csname peffConception\endcsname{2}
\expandafter\def\csname peffRéalisation\endcsname{1}
\expandafter\def\csname peffExploitation\endcsname{2}
\expandafter\def\csname rimpR01\endcsname{4}
\expandafter\def\csname rvraR01\endcsname{3}
\expandafter\def\csname rscoreR01\endcsname{12}
\expandafter\def\csname rimpR02\endcsname{4}
\expandafter\def\csname rvraR02\endcsname{3}
\expandafter\def\csname rscoreR02\endcsname{12}
\expandafter\def\csname rimpR03\endcsname{4}
\expandafter\def\csname rvraR03\endcsname{2}
\expandafter\def\csname rscoreR03\endcsname{8}
\expandafter\def\csname rimpR04\endcsname{3}
\expandafter\def\csname rvraR04\endcsname{3}
\expandafter\def\csname rscoreR04\endcsname{9}
\expandafter\def\csname rimpR05\endcsname{2}
\expandafter\def\csname rvraR05\endcsname{2}
\expandafter\def\csname rscoreR05\endcsname{4}
\expandafter\def\csname rimpR06\endcsname{3}
\expandafter\def\csname rvraR06\endcsname{3}
\expandafter\def\csname rscoreR06\endcsname{9}
\expandafter\def\csname acoutA01\endcsname{18\,000\,€}
\expandafter\def\csname acoutA02\endcsname{16\,000\,€}
\expandafter\def\csname acoutA03\endcsname{14\,000\,€}
\expandafter\def\csname acoutA04\endcsname{24\,000\,€}
\expandafter\def\csname acoutA05\endcsname{22\,000\,€}
\expandafter\def\csname acoutA06\endcsname{35\,000\,€}
\expandafter\def\csname acoutA07\endcsname{14\,000\,€}
\expandafter\def\csname acoutA08\endcsname{28\,000\,€}
\expandafter\def\csname acoutA09\endcsname{18\,000\,€}
\expandafter\def\csname acoutA10\endcsname{30\,000\,€}
\expandafter\def\csname acoutA11\endcsname{35\,000\,€}


\definecolor{encre}{HTML}{1F2A37}\definecolor{accent}{HTML}{1D5C7A}\definecolor{accentclair}{HTML}{E4EEF3}
\definecolor{alerte}{HTML}{A33A3A}\definecolor{vert}{HTML}{3B7A57}\definecolor{gris}{HTML}{5F6B7A}\definecolor{grisclair}{HTML}{F2F4F6}
\setbeamercolor{normal text}{fg=encre}\setbeamercolor{frametitle}{fg=accent}\setbeamercolor{structure}{fg=accent}
\setbeamerfont{frametitle}{size=\Large,series=\bfseries}
\setbeamertemplate{navigation symbols}{}
\setbeamertemplate{footline}{\hfill{\scriptsize\color{gris}MEDICA-SERVICES · CODIR · \insertframenumber/\inserttotalframenumber}\hspace*{1em}\vspace{0.6em}}
\newcolumntype{Y}{>{\raggedright\arraybackslash}X}

\newcommand{\troiscartes}[9]{% couleur, titre, texte pour trois cartes alignées
  \begin{tikzpicture}[carte/.style={rounded corners=3pt, text width=4.0cm, inner sep=7pt, align=left, anchor=north}]
    \node[carte, draw=#1, fill=#1!6] at (0,0) {\parbox[t][4.9cm][t]{\linewidth}{{\bfseries\color{#1}#2}\\[4pt]\small #3}};
    \node[carte, draw=#4, fill=#4!6] at (4.65,0) {\parbox[t][4.9cm][t]{\linewidth}{{\bfseries\color{#4}#5}\\[4pt]\small #6}};
    \node[carte, draw=#7, fill=#7!6] at (9.3,0) {\parbox[t][4.9cm][t]{\linewidth}{{\bfseries\color{#7}#8}\\[4pt]\small #9}};
  \end{tikzpicture}}

\begin{document}

\begin{frame}[plain]
  \begin{tikzpicture}[remember picture,overlay]
    \fill[accent] (current page.north west) rectangle ([yshift=-5cm]current page.north east);
    \node[anchor=north west, text=white, align=left, text width=13cm] at ([xshift=1.2cm,yshift=-1.2cm]current page.north west) {%
      {\footnotesize\Formation}\\[8pt]{\huge\bfseries Audit de maturité SSI}\\[5pt]{\Large MEDICA-SERVICES · Restitution au CODIR}};
    \node[anchor=north west, align=left, text width=13.5cm] at ([xshift=1.2cm,yshift=-5.7cm]current page.north west) {%
      {\small \Membres\ · \DateDepot}\\[6pt]
      {\footnotesize\color{gris} Audit sur pièces P01 à P17 et T01 à T10. Méthode SGDN/DCSSI 2007 utilisée comme référence pédagogique, sans valeur de certification. Cas fictif.}};
  \end{tikzpicture}
\end{frame}

\begin{frame}{Où en est MEDICA-SERVICES}
  \begin{columns}[T]
  \begin{column}{0.46\textwidth}
  \textbf{Niveau adéquat : \Cible\ sur 5}\par\smallskip
  {\small Somme des maxima \SommeMaxima. Données de santé, sécurité de patients chez 12 clients, reprise contractuelle en 8 h.}\par\bigskip
  \textbf{Niveau démontré : 1 à 2}\par\smallskip
  {\small Aucun processus au niveau 3. Écart de 2 à 3 niveaux.}\par\bigskip
  {\small\color{alerte} Une affirmation contredite (MFA) et une non démontrée (tests de restauration), P17.}
  \end{column}
  \begin{column}{0.5\textwidth}
  \small
  \begin{tabularx}{\linewidth}{@{}Y c c@{}}
  \toprule Processus & Démontré & Cible \\ \midrule
  Stratégie & \peff{Stratégie} & \Cible \\ Risques & \peff{Risques} & \Cible \\ Règles & \peff{Règles} & \Cible \\
  Supervision & \peff{Supervision} & \Cible \\ Conception & \peff{Conception} & \Cible \\ Réalisation & \peff{Réalisation} & \Cible \\
  Exploitation & \peff{Exploitation} & \Cible \\ \bottomrule
  \end{tabularx}
  \end{column}
  \end{columns}
\end{frame}

\begin{frame}{Trois risques majeurs}
  \centering
  \troiscartes{alerte}{R01 · Compte à privilèges détourné}{Score \rscore{R01}. 4 comptes privilégiés sur 10 non enrôlés MFA, dont un administrateur de domaine sans activité documentée depuis 2025 (T03). Modification non autorisée de l'AD et des serveurs.}{alerte}{R02 · Reprise au-delà de 8 h}{Score \rscore{R02}. Rançongiciel ou panne. Restaurations peu éprouvées et sauvegardes joignables depuis tous les postes.}{accent}{R03 · Données patients via APP01}{Score \rscore{R03}. Correctif critique manquant, hôte en retard de 74 jours, sans exception. Impact maximal, levier rapide. Application non exposée directement.}
\end{frame}

\begin{frame}{Trois constats décisifs}
  \centering
  \troiscartes{alerte}{Alerte AL-1048 sans examen tracé}{9 octobre, 00:13 à 00:22 UTC. Compte admin-old connecté au VPN après 3 échecs, ajout à Server-Admins. Fermeture automatique 24 h plus tard, aucun ticket. À investiguer, pas une intrusion prouvée (T07, T08).}{alerte}{Restauration testée sur 2 services sur 9}{La DG déclare des tests trimestriels de toutes les restaurations. Deux PV cités, contenu non communiqué. RTO de 8 h non démontré, RPO non fourni (T09, P17).}{accent}{MFA à 60 \% des comptes privilégiés}{La DSI déclare 100 \%. Export : 6 comptes enrôlés sur 10. Prestataire et compte de sauvegarde sans second facteur (T03, T04, P17).}
\end{frame}

\begin{frame}{Trois décisions demandées}
  \small
  \begin{tabularx}{\linewidth}{@{}c Y Y@{}}
  \toprule & Décision & Effet attendu \\ \midrule
  \textbf{1} & Mandater dès aujourd'hui l'investigation d'AL-1048 et les mesures conservatoires (admin-old, FW-31, FW-17) & Événement qualifié, accès à risque fermés sous un mois \\[3pt]
  \textbf{2} & Engager \TotalRetenu\ sur \Plafond\ pour six actions (A01, A02, A03, A05, A08, A09) & Les 7 processus au niveau 2 démontré à 12 mois, marge \Marge \\[3pt]
  \textbf{3} & Accepter par écrit pour 12 mois les risques résiduels des mesures différées, dont la segmentation (A10) et la standardisation (A04) & R04, R06 et part de R02 suivis et réexaminés au bilan annuel \\ \bottomrule
  \end{tabularx}
\end{frame}

\begin{frame}{Feuille de route sur 12 mois}
  \begin{tikzpicture}[x=0.72cm,y=0.55cm,font=\small]
    \foreach \m in {1,...,12}{\node[gris,font=\scriptsize] at (\m-0.5,0.8) {M\m};}
    \foreach \l/\a/\b/\t [count=\i] in {Suites de l'investigation/0/1/{Clôture de l'incident AL-1048},
      A08 MFA et privilèges/0/3/100 \% MFA ou exception,A02 Supervision/0/3/Alertes tracées,
      A01 Risques/1/4/Registre approuvé,A09 Correctifs/0/4/{APP01 en M1, SLA en M4},A05 Restauration/2/6/9 PV sur 9,
      A03 PSSI et sensibilisation/3/6/PSSI v2,Stabilisation/6/12/Bilan de maturité}{
      \fill[vert!80] (\a,-\i+0.3) rectangle (\b,-\i-0.3);
      \node[anchor=east,font=\scriptsize] at (-0.1,-\i) {\l};
      \node[anchor=west,font=\scriptsize,gris] at (\b+0.05,-\i) {\t};}
  \end{tikzpicture}
  \par\smallskip{\footnotesize\color{gris} Hypothèse : décision du CODIR à la mi-novembre 2026, M1 = décembre 2026. Mesures conservatoires engagées dès la remise du rapport. Chaque palier est validé par une preuve de clôture, non par déclaration.}
\end{frame}

\end{document}
